\begin{workedexample}[Worked Example: A gain chain in dBm]
An antenna delivers $-40$~dBm to a front end with a $+15$~dB LNA, a
$-6$~dB filter, and $+30$~dB of IF gain. Because gains multiply as powers but
\emph{add} as decibels, the output level is
\[ P_{out} = -40 + 15 - 6 + 30 = -1\ \text{dBm}, \]
i.e.\ $P = 10^{-0.1} = 0.79\ \text{mW}$. Working in dB turns a product of linear
power ratios into a running sum.
\end{workedexample}

\begin{keytakeaways}
\begin{itemize}[leftmargin=1.4em,itemsep=2pt]
\item $P_{\text{dBm}} = 10\log_{10}(P/1\,\text{mW})$; $0$~dBm $=1$~mW, $+30$~dBm $=1$~W.
\item Cascaded gains and losses simply add in dB.
\item dBc is relative to a carrier, dBi is antenna gain vs.\ isotropic, dBW $=$ dBm $-30$.
\end{itemize}
\end{keytakeaways}

\begin{exercises}
\item Convert $+43$~dBm to watts.
\item A signal is $-70$~dBm and a spur is $-90$~dBm. What is the spur level in dBc?
\item Express $+30$~dBm in dBW.
\end{exercises}

\furtherreading{Pozar~\cite{pozar} (ch.~10); Hayward~\cite{hayward}.}
